% Entanglement chapter.

\section{Distributional entanglement}\label{sec:entanglement}

The witness pair of Section~\ref{sec:warmup} is a specific pair of six-vertex
graphs. On a reading of the paper by an applied audience, the natural
question is whether the impossibility is a curiosity of small hand-crafted
graphs, or whether it survives on distributions of graphs of arbitrary
size. This chapter answers by lifting the impossibility from the pair to
a family of ensemble pairs indexed by graph size. Every sample from
either ensemble in the family carries the properties of the base
witness: the two ensembles agree on the $1$-WL colour multiset, and the
clustering coefficient differs deterministically by $1$.

\subsection{Ensemble notation}

An \emph{ensemble} of featured graphs on $n$ vertices is a probability
distribution over the set of featured graphs on $n$ vertices with
features in a fixed set $\mathcal{X}$. Write $\mathrm{supp}(\mathcal{M})$
for the support of the distribution $\mathcal{M}$: the set of featured
graphs the distribution assigns positive mass to. The results below
concern the support pointwise, not the distribution over the support,
so no measure-theoretic subtleties arise; sampling only ensures that
every graph in the support is a plausible instance.

\subsection{The two ensembles}\label{sec:two-ensembles}

Fix $n$ divisible by $12$ and features in $\{-1, +1\}$.

Define $\mathcal{M}_n^A$ as the following ensemble. Draw an ordered
partition of $[n]$ into $n/3$ triples uniformly at random. The first
$n/6$ triples are the \emph{$(+1)$-triangles}; the remaining $n/6$ are
the \emph{$(-1)$-triangles}. The graph $\gph_A$ consists of the
disjoint union of triangles, one triangle on each triple. The feature
$x_A(v)$ is $+1$ if $v$ belongs to a $(+1)$-triangle and $-1$
otherwise.

Define $\mathcal{M}_n^B$ analogously with $6$-cycles in place of
triangles. Draw an ordered partition of $[n]$ into $n/6$ ordered
$6$-tuples uniformly at random. The first $n/12$ tuples are
\emph{$(+1)$-cycles}, the remaining $n/12$ are \emph{$(-1)$-cycles}.
The graph $\gph_B$ consists of the disjoint union of $C_6$s, one on
each tuple (with the cyclic edge structure). The feature $x_B(v)$ is
$+1$ if $v$ belongs to a $(+1)$-cycle and $-1$ otherwise.

Every featured graph in $\mathrm{supp}(\mathcal{M}_n^A) \cup
\mathrm{supp}(\mathcal{M}_n^B)$ has $n$ vertices, every vertex has
degree $2$, exactly $n/2$ vertices carry feature $+1$, and exactly
$n/2$ carry $-1$. The two ensembles differ in their component
structure: samples from $\mathcal{M}_n^A$ decompose into $n/3$
triangles; samples from $\mathcal{M}_n^B$ decompose into $n/6$
$6$-cycles.

\subsection{Fingerprint equality and closure gap}

The two ensembles satisfy the fingerprint and clustering conditions
required by Theorem~\ref{thm:main}. Setting $\mathcal{A}_n =
\mathrm{supp}(\mathcal{M}_n^A)$ and $\mathcal{B}_n =
\mathrm{supp}(\mathcal{M}_n^B)$, we prove them here.

\begin{proposition}\label{prop:entanglement}
Fix $n \equiv 0 \pmod{12}$. For every $(\gph_A, x_A) \in
\mathcal{A}_n$ and $(\gph_B, x_B) \in \mathcal{B}_n$,
\[
\mathrm{hist}(\gph_A, x_A) = \mathrm{hist}(\gph_B, x_B),
\qquad \Pi_{\clos}(\gph_A) = 1, \qquad \Pi_{\clos}(\gph_B) = 0.
\]
\end{proposition}

\begin{proof}
\emph{Fingerprint equality.} Fix a sample $(\gph_A, x_A) \in
\mathcal{A}_n$. Consider $1$-WL refinement on $(\gph_A, x_A)$.

At round $0$, colours are $c_v^{(0)} = \tau(x_v) \in \{\tau(+1),
\tau(-1)\}$. Call these two colours $\alpha_+$ and $\alpha_-$. There
are $n/2$ vertices with colour $\alpha_+$ and $n/2$ with colour
$\alpha_-$.

At round $1$, every vertex $v \in (+1)$-triangle has two neighbours,
both in the same triangle, so both carry feature $+1$. The refined
colour is $c_v^{(1)} = \tau(\alpha_+, \{\!\{\alpha_+, \alpha_+\}\!\}) =
\beta_+$. Analogously, every vertex in a $(-1)$-triangle receives
colour $\beta_-$.

At round $2$, every vertex has two neighbours, both coloured
$\beta_+$ (in a $(+1)$-triangle) or both coloured $\beta_-$ (in a
$(-1)$-triangle). The refined colour is $\gamma_+$ or $\gamma_-$. This
is the same partition of vertices as at round $1$, so refinement has
stabilised. The graph-level fingerprint is
\[
\mathrm{hist}(\gph_A, x_A) = \{\!\{\gamma_+^{n/2},\ \gamma_-^{n/2}\}\!\}.
\]

Fix a sample $(\gph_B, x_B) \in \mathrm{supp}(\mathcal{M}_n^B)$. In a
$(+1)$-cycle, every vertex has two neighbours in the same cycle,
which by construction carry the same feature $+1$. Round-one $1$-WL
gives $c_v^{(1)} = \beta_+$ at every such vertex. Round-two and
beyond stabilise at $\gamma_+$. Analogously for $(-1)$-cycles.
Hence
\[
\mathrm{hist}(\gph_B, x_B) = \{\!\{\gamma_+^{n/2},\ \gamma_-^{n/2}\}\!\} =
\mathrm{hist}(\gph_A, x_A).
\]
The two featured graphs are $1$-WL equivalent.

\emph{Closure gap.} In a triangle, each vertex has degree $2$ and
its two neighbours are joined by an edge (being the other two
vertices of the triangle), so the local clustering coefficient of
every vertex in a $(+1)$-triangle or $(-1)$-triangle is $1$.
Averaged over $\gph_A$, $\Pi_{\clos}(\gph_A) = 1$. In a $6$-cycle,
each vertex has degree $2$ and its two neighbours are at cyclic
distance $2$ apart, not joined, so the local clustering coefficient
is $0$ at every vertex. Averaged over $\gph_B$, $\Pi_{\clos}(\gph_B)
= 0$.
\end{proof}

Proposition~\ref{prop:entanglement} discharges the two hypotheses of
Theorem~\ref{thm:main}. Item (i) of the main theorem follows from
$\mathrm{hist}(\gph_A, x_A) = \mathrm{hist}(\gph_B, x_B)$ combined
with Proposition~\ref{prop:upper} (any MPNN with continuous update and readout
is $1$-WL invariant). Items (ii), (iii), (iv) of the main theorem
follow from the closure statement above and the MSE calculation in
Section~\ref{sec:main}.

\begin{corollary}\label{cor:no-closure-mpnn}
No MPNN estimator can approximate any polarisation measure in the
right column of the taxonomy of Section~\ref{sec:measures} uniformly on
any class $\mathcal{G}$ that includes samples from both
$\mathcal{A}_n$ and $\mathcal{B}_n$ for some $n$.
\end{corollary}

\begin{proof}
The listed measures separate every pair $(\gph_A, \gph_B) \in
\mathcal{A}_n \times \mathcal{B}_n$ by an amount bounded away from
zero: $\Pi_{\clos}$ and $t$ are covered by the closure statement of
Proposition~\ref{prop:entanglement}; $\mathrm{GED}$ is covered by the
per-component witness of Proposition~\ref{prop:noncomputable-measures}. By $\mathrm{hist}(\gph_A, x_A) =
\mathrm{hist}(\gph_B, x_B)$ and Proposition~\ref{prop:upper}, every MPNN gives
identical outputs across every such pair. Uniform approximation of
a functional that separates the two families by a constant is
impossible.
\end{proof}

\subsection{Why the ensembles are structurally uniform}

Each connected component of a sample from $\mathcal{M}_n^A$ or
$\mathcal{M}_n^B$ is a triangle or a $6$-cycle, and each component is
monochromatic in the feature. This is what makes the round-one $1$-WL
computation clean: every vertex has neighbour features that agree with
its own feature, and no cross-feature edges exist. The absence of
cross-feature edges is what puts both ensembles on the same $1$-WL
fingerprint.

If one weakens the ensembles by adding cross-feature edges or by
letting features be assigned independently of component structure, the
round-one $1$-WL colour of a vertex depends on its neighbourhood
feature composition, and matching the two ensembles at round one
requires matching the joint distribution of (feature, neighbour
feature multiset) tuples. The next subsection records the extent to
which one can weaken the construction and preserve the impossibility.

\subsection{Attempts at ensembles that resemble real
interaction graphs}\label{sec:obstacle}

The ensembles above are structurally uniform: every connected
component is a triangle or a $6$-cycle, and each component is
monochromatic. Real interaction data does not decompose this way. It
is natural to ask whether the impossibility survives on ensembles
that better model real data. We attempted five weakenings; each falls
short in a specific measurable way. The attempts are documented here
because the pattern is informative about the scope of the
impossibility, and because a reader hoping to prove a stronger
theorem in this line will need to move around each obstacle.

\subsubsection*{Two-block planted-partition model}

Fix $n$, integers $p, q \in [0, 1]$, and a real $\sigma \in [0, 1/2]$.
The \emph{feature-perturbed planted-partition model} $\ppm(p, q,
\sigma)$ on $n$ vertices proceeds as follows.

Partition $[n]$ evenly into two blocks $B_+$ and $B_-$ of size
$n/2$. For each pair $\{u, v\}$ of vertices, place an edge
independently with probability $p$ if $u$ and $v$ belong to the same
block, and $q$ otherwise. Assign an intermediate feature $\tilde x_v$
of $+1$ if $v \in B_+$ and $-1$ if $v \in B_-$. Then independently
flip each $\tilde x_v$ with probability $\sigma$ to obtain the
observed feature $x_v$.

The model has three parameters. Matching two of them (expected degree
and edge homophily) between two parameter points leaves one degree of
freedom, along which the expected clustering coefficient varies. A
grid sweep at $n = 1000$ confirms that on the constant-degree-and-homophily
level set of parameter space, the expected clustering ranges over an
interval of width at least $0.07$.

The obstruction is at $1$-WL round one. Matching a scalar homophily
statistic does not match the depth-$1$ joint distribution of $(x_v,
+\text{count}, -\text{count})$ tuples over $v$. Numerical experiment:
between two matched $\ppm(p, q, \sigma)$ points with common expected
degree and common edge homophily, the depth-$1$ $1$-WL type
distributions have total variation distance $0.28$. The $\ppm$
ensembles are not $1$-WL indistinguishable to depth $1$, and the
matching cannot be improved with three free parameters.

\subsubsection*{$K$-block stochastic block model}

Adding blocks adds parameters. A stochastic block model with $K$
blocks per feature side, block size $m = n/(2K)$, and edge
probabilities $p_{\mathrm{in}}$ (within-block), $p_{\mathrm{sfd}}$
(same-feature different-block), $p_{\mathrm{df}}$ (different-feature)
was tested at $K = 4$ with a swap construction: parameter point $1$
has $(p_{\mathrm{in}}, p_{\mathrm{sfd}}) = (0.18, 0.02)$, parameter
point $2$ has $(0.02, 0.18)$; both keep the other parameters equal.
The two points have the same first-moment neighbour-count
distribution for both feature classes. Numerical experiment: depth-$1$
$1$-WL total variation between the two ensembles is $0.07$, and the
clustering gap is $0.045$.

A numerical solver matching the first \emph{two} moments of the
neighbour-count distribution at $K \in \{3, 4, 5, 6, 8, 10\}$ per
feature side, subject to $K \cdot m$ constant, reduces the
depth-$1$ $1$-WL total variation to the range $[0.06, 0.09]$ and
shrinks the clustering gap to the range $[0.01, 0.02]$. The ratio
between the two does not improve.

The obstruction: matching the full joint distribution of $T$-hop
neighbour compositions between two parameter points of an SBM
requires matching every moment of the neighbour-count distribution.
A finite-parameter SBM has finitely many moments to spare;
simultaneously leaving the clustering coefficient free consumes them
faster than adding blocks supplies them.

\subsubsection*{Random $d$-regular graph and its two disjoint halves}

Take two ensembles on $n$ vertices at fixed degree $d$: (i) a random
$d$-regular graph on $[n]$, and (ii) a graph consisting of two
disjoint random $d$-regular graphs, each on $n/2$ vertices. Both
ensembles have the same Benjamini-Schramm local weak limit as $n
\to \infty$, namely the infinite $d$-regular tree.

Under an opinion vector $o$ aligned with the natural bipartition of
the disconnected ensemble (so $o$ is $+1$ on one component of (ii)
and $-1$ on the other), the featured graphs are $1$-WL
distinguishable at round one. In the disconnected ensemble every
neighbour of a $+1$-vertex is $+1$, so the depth-$1$ colour is
$\tau(+1, \{\!\{+1\}\!\}^d)$. In the connected ensemble, half of a
$+1$-vertex's $d$ neighbours are $+1$ (in expectation), so the
depth-$1$ colour distribution is spread across every value of
$(\text{same-feature count}, \text{different-feature count})$.
Numerical experiment: depth-$1$ $1$-WL total variation between the
two featured ensembles is $0.94$. Removing the opinion vector
restores $1$-WL equivalence and removes the polarisation signal.

\subsubsection*{Cai-F\"{u}rer-Immerman construction}

The Cai-F\"{u}rer-Immerman construction, developed for the graph
isomorphism problem \citep{CFI1992}, produces pairs of graphs $(G^0,
G^1)$ that are non-isomorphic but agree on every $k$-WL colouring
for a specified $k$. The construction operates by attaching a
CFI-gadget to each vertex of a base graph; the choice of gadget
determines whether $G^0$ or $G^1$ is obtained.

For $k = 1$, the CFI construction gives pairs indistinguishable by
$1$-WL. However, the construction preserves every local homomorphism
count up to graphs of size $k + 1 = 2$, which is trivial; more
relevantly, the standard CFI gadgets preserve triangle counts. The
construction therefore cannot generate a closure gap between the two
sides of the pair.

\subsubsection*{Graphon limit}

For ensembles converging in the graphon sense to a graphon $W : [0,
1]^2 \to [0, 1]$, message-passing GNNs are universal approximators
of continuous permutation-invariant graphon functionals under the
cut metric \citep{KerivenPeyre2020, RuizChamonRibeiro2021}. Triangle
density on such ensembles converges to the graphon homomorphism
density $t(K_3, W) = \int W(x_1, x_2) W(x_2, x_3) W(x_3, x_1) \,
dx_1 dx_2 dx_3$, and this quantity is Lipschitz in the cut metric
on $W$ \citep{LovaszBook}.

In the dense graphon limit, closure is asymptotically MPNN-computable.
A stronger ensemble-level impossibility on random-graph ensembles
that resemble real interaction data has to operate in the sparse or
structured regime rather than the dense graphon limit.

\subsection{Summary of the obstacles}

The above five attempts sketch the shape of what is and is not
available. The impossibility survives on structured ensembles of
arbitrary size (Proposition~\ref{prop:entanglement}) and on ensembles with
suitable degree and homophily matching (the PPM sweep), but the
$1$-WL distance between the two matched sides shrinks slower than
the closure gap as more matching constraints are added. In the dense
graphon limit, the impossibility disappears entirely.

The empirical audit of Section~\ref{sec:emp-reddit} runs on Twitter
retweet graphs, which sit between the structured ensembles of
Proposition~\ref{prop:entanglement} and the dense graphon regime. The audit
tests directly whether a published MPNN estimator is invariant to
closure-destroying rewiring on such graphs; the answer, as reported
in Section~\ref{sec:emp-results}, is yes.
